% =====================================================================
%  CHAPTER 3: 儿童少年身体发育（对应大纲第四章）
%  参考往届笔记：第三章 儿童少年身体发育
% =====================================================================
\chapter{儿童少年身体发育}
\setcounter{chapter}{3}
\chapterintro{
\textbf{【掌握】}儿童少年体格发育（体格阶段性变化、身体比例变化、体型变化）；儿童少年体能发育（健康相关体能、运动技能相关体能、体能发育特点、体力活动）。\par
\textbf{【熟悉】}儿童少年体成分发育（体成分模型、体成分发育的性别特征）。
}

% ----- 5.1 体格发育 -----
\section{儿童少年体格发育}

\subsection{体格阶段性变化}

\begin{keybox}
\textbf{体格发育（Physical Growth）：}人体外部形态、身体比例和体型等方面随年龄发生的变化。

\textbf{一、体格阶段性变化：}身高$\to$身体线性生长（Linear Growth），体重$\to$重量。

\textbf{1. 第一生长突增期：}胎儿$\to$2岁。
\begin{itemize}[leftmargin=*]
    \item \textbf{身高：}孕中期$\to$1岁末；孕中期（4-6个月）增长量占胎儿期总增长量的50\%；出生后第1年增长约25cm。
    \item \textbf{体重：}孕晚期$\to$1岁末；孕晚期（7-9个月）增长量占胎儿期总增长量的70\%。
\end{itemize}

\textbf{2. 相对稳定期：}2岁$\to$青春期发育前。身高：$\uparrow$5-7cm/年，体重：$\uparrow$2-3kg/年。

\textbf{3. 第二生长突增期/青春期生长突增（Adolescent Growth Spurt）：}青春期。
\begin{itemize}[leftmargin=*]
    \item \textbf{身高：}突增高峰年增长10-14cm，男童$>$女童。
    \item \textbf{体重：}每年增长4-7kg，突增高峰年增长8-10kg。
\end{itemize}

\textbf{4. 生长停滞期：}青春后期开始。
\end{keybox}

\subsection{身体比例变化}

\begin{keybox}
\textbf{二、身体比例变化}

\textbf{1. 反映横向身体比例指标：}反映身体各横截面指标的相互关系。
\begin{itemize}[leftmargin=*]
    \item \imp{身高胸围指数}（胸围/身高）和\imp{肩盆宽指数}（肩宽/盆宽）：只能衡量局部（身体上部或下部），实测值重复性差，后者个体差异只在青春期显现。
    \item \imp{BMI}：既能反映身体胖瘦，又能确切反映不同群体、年代身体充实度变化。许多发达国家特别重视对BMI高百分位数变化的动态监测，将其作为制定肥胖防治策略和措施的重要依据。
\end{itemize}

\textbf{2. 反映纵向身体比例指标：}反映个体下肢长度和线性高度（身高）的比例关系。
\begin{itemize}[leftmargin=*]
    \item \imp{身高坐高指数} = 坐高/身高。
    \item \imp{下肢长指数：}下肢长指数I = (身高$-$坐高)/身高 $\times$ 100\%，下肢长指数II = (身高$-$坐高)/坐高 $\times$ 100\%。
\end{itemize}
\end{keybox}

\subsection{体型变化}

\begin{defbox}
\textbf{三、体型变化}

\textbf{1. 体型（Somatotype）：}对人体形态的总体描述和评定，是身体各部位大小比例的形态特征总和，由遗传因素决定，也受人体对环境的适应和诸多环境因素的综合影响。

\textbf{2. 分类}
\begin{itemize}[leftmargin=*]
    \item \imp{内胚层体型（Endomorphy）}：身体圆胖，消化器官发达。
    \item \imp{中胚层体型（Mesomorphy）}：身体健壮，骨骼肌肉发达。
    \item \imp{外胚层体型（Ectomorphy）}：身体瘦长，神经感觉系统发达。
\end{itemize}
\end{defbox}

% ----- 5.2 体能发育 -----
\section{儿童少年体能发育}

\subsection{概念}

\begin{keybox}
\textbf{体能（Physical Fitness）/ 体适能：}人体具备的能胜任日常工作和学习而不感到疲劳，同时有余力能充分享受休闲娱乐生活，又可应付突发紧急情况的能力。

\textbf{一、概念}

\textbf{1. 健康相关体能（Health-Related Physical Fitness）：}体能的基础部分，其目标是为身体健康和高质量，泛指人的体质，维持身体健康、提高工作、学习和生活效率所必需的基本技能。
\begin{itemize}[leftmargin=*]
    \item 指标体系：体成分、心肺耐力、柔韧性、肌力、肌耐力。
    \item 是儿童少年为维持健康、提高生活质量等需求的追求。
\end{itemize}

\textbf{2. 运动技能相关体能（Sports-Related Physical Fitness）：}建立在健康相关体能基础上，属于较高体能需求层次，目标是为比赛取胜奖项。
\end{keybox}

\subsection{儿童少年体能发育特点}

\begin{keybox}
\textbf{二、儿童少年体能发育特点}

\textbf{1. 不均衡性：}不同体能指标在不同年龄发育速度有快有慢。

\textbf{2. 阶段性}
\begin{itemize}[leftmargin=*]
    \item 快速增长阶段：男童6-14岁，女童6-12岁。
    \item 慢速增长阶段：男童15-18岁，女童12-15岁。
    \item 恢复性生长：仅女童16-18岁。
    \item 稳定阶段：男性19-25岁，女性19-22岁。
\end{itemize}

\textbf{3. 不平衡性}

\textbf{4. 性别特征}
\end{keybox}

\subsection{儿童少年体力活动}

\begin{keybox}
\textbf{三、儿童少年体力活动}

\textbf{体力活动（Physical Activity）：}WHO定义，由骨骼肌收缩产生需要能量消耗的任何身体运动，包括进行工作、家务、旅游、游戏、娱乐等。
\end{keybox}

% ----- 5.3 体成分发育 -----
\section{儿童少年体成分发育}

\subsection{体成分模型}

\begin{defbox}
\textbf{身体成分（Body Composition）/ 体成分：}人体总重量中不同身体成分的构成比例，属化学生长（Chemical Growth）范畴。

\textbf{一、体成分模型}

\textbf{1. 二成分模型：}体脂重（FM）+ 去脂体重（FFM），主要使用。

\textbf{2. 多成分模型：}二成分模型的细化。
\end{defbox}

\subsection{体成分发育的性别特征}

\begin{defbox}
\textbf{二、体成分发育的性别特征}

\textbf{1. 体成分发育的性别差异}
\begin{itemize}[leftmargin=*]
    \item 生命早期，男女体脂量接近。
    \item 儿童期，男、女体脂均下降。
    \item 青春期，女童变化以体脂增加为主，男童体脂有下降。
    \item 体脂分布的性别差异在青春期前即有表现；男童上半身的皮下和内脏脂肪蓄积，女童皮下脂肪蓄积在臀部。
\end{itemize}

\textbf{2. 瘦体重发育的性别特征}
\begin{itemize}[leftmargin=*]
    \item 生命早期开始，男孩就显出更高的瘦体重，故其体重$>$同年龄女孩。
    \item 儿童期，男孩仍较女孩有较多瘦体重增长。
    \item 进入青春期后，男孩主要表现为瘦体重增加，女孩的瘦体重增量仅为男孩的一半。
    \item 随着近年来性早熟（女性更多）现象增加，瘦体重的性别差异出现时间明显提前。
\end{itemize}
\end{defbox}

\newpage
